% Part: first-order-logic
% Chapter: tableaux
% Section: soundness.tex

\documentclass[../../../include/open-logic-section]{subfiles}

\begin{document}

\iftag{FOL}
      {\olfileid[fr]{fol}{tab}{sou}}
      {\olfileid[fr]{pl}{tab}{sou}}
      
\olsection{Correction}

\begin{explain}
Un système de !!{derivation}, tel que celui des tableaux, est
\emph{correct} s'il ne permet pas de !!{derive} des résultats qui ne
valent pas réellement. La correction constitue ainsi une garantie de
fiabilité pour les systèmes de !!{derivation}. Selon la propriété de
théorie de la démonstration considérée, nous voulons notamment savoir
que :
\begin{enumerate}
\item tout $!A$ !!{derivable} est \iftag{FOL}{valide}{une tautologie};
\item si !!a{sentence} est !!{derivable} à partir d'autres, il en est
  aussi une conséquence;
\item si un ensemble d'!!{sentence}s est incohérent, il est
  insatisfaisable.
\end{enumerate}
Ces propriétés sont essentielles pour un système de !!{derivation}.
Si l'une d'elles était en défaut, le système serait défectueux : il
permettrait de !!{derive} trop de résultats. Il est donc primordial
d'en établir la correction.

Comme toutes ces propriétés de théorie de la démonstration sont
définies au moyen de !!{tableau}s fermés, la preuve de (1)--(3) exige
une propriété sémantique de ces tableaux. Nous définirons d'abord ce
que signifie, pour une \iftag{FOL}{structure}{valuation}, satisfaire
!!a{signed formula}, puis nous montrerons que, si !!a{tableau} est
fermé, aucune \iftag{FOL}{structure}{valuation} ne satisfait toutes ses
hypothèses. Les propriétés (1)--(3) en découleront comme corollaires.
\end{explain}

\begin{defn}
\iftag{FOL}{!!^a{structure}}{!!^a{valuation}}~$\iftag{FOL}{\Struct{M}}{\pAssign{v}}$
\emph{satisfait} !!a{signed formula} $\sFmla{\True}{!A}$ si et seulement si
$\iftag{FOL}{\Sat{M}}{\pSat{v}}{!A}$, et satisfait
$\sFmla{\False}{!A}$ si et seulement si
$\iftag{FOL}{\Sat/{M}}{\pSat/{v}}{!A}$. $\iftag{FOL}{\Struct{M}}{\pAssign{v}}$
satisfait un ensemble d'!!{signed formula}s~$\Gamma$ si et seulement
s'il satisfait chaque $\sFmla{S}{!A} \in \Gamma$. L'ensemble
$\Gamma$ est \emph{satisfaisable} s'il existe
\iftag{FOL}{!!a{structure}}{!!a{valuation}} qui le satisfait, et
\emph{insatisfaisable} dans le cas contraire.
\end{defn}

\begin{thm}[Correction]
  \ollabel{thm:tableau-soundness}
  Si $\Gamma$ possède un !!{tableau} fermé, alors $\Gamma$ est
  insatisfaisable.
\end{thm}

\begin{proof}
Disons qu'une branche d'!!a{tableau} est satisfaisable si et seulement
si l'ensemble des !!{signed formula}s qui y figurent est satisfaisable,
et qu'!!a{tableau} est satisfaisable s'il contient au moins une branche
satisfaisable.

Montrons que toute extension d'!!a{tableau} satisfaisable par une règle
d'inférence donne encore un !!{tableau} satisfaisable. Le théorème en
découlera : tout !!{tableau} fermé s'obtient en appliquant des règles
d'inférence au !!{tableau} qui ne contient que les hypothèses de
$\Gamma$. Si $\Gamma$ était satisfaisable, tout !!{tableau} pour
$\Gamma$ le serait donc aussi. Or un !!{tableau} fermé ne peut être
satisfaisable : chaque branche contient à la fois
$\sFmla{\True}{!A}$ et $\sFmla{\False}{!A}$, et aucune
\iftag{FOL}{structure}{valuation} ne peut à la fois satisfaire et ne
pas satisfaire~$!A$.

Supposons donné un !!{tableau} satisfaisable, c'est-à-dire un
!!{tableau} qui possède au moins une branche satisfaisable. L'application
d'une règle d'inférence ajoute des !!{signed formula}s à une branche,
ou bien ramifie une branche en deux. Si le !!{tableau} possède une
branche satisfaisable qui n'est pas prolongée par l'application
considérée, cette branche demeure satisfaisable dans le !!{tableau}
prolongé; celui-ci est donc satisfaisable. Il suffit par conséquent
d'examiner le cas où la règle est appliquée à une branche satisfaisable.

Soit $\Gamma$ l'ensemble des !!{signed formula}s de cette branche, et
soit $\sFmla{S}{!A} \in \Gamma$ la !!{signed formula} à laquelle la
règle est appliquée. Si la règle ne provoque pas de ramification, il
faut montrer que la branche prolongée, c'est-à-dire $\Gamma$ avec les
conclusions de la règle, demeure satisfaisable. Si la règle provoque
une ramification, il faut montrer qu'au moins l'une des deux branches
obtenues est satisfaisable.
  
Considérons d'abord les inférences qui ne provoquent pas de ramification.
\begin{enumerate}
\item La branche est prolongée par l'application de
  $\TRule{\True}{\lnot}$ à
  $\sFmla{\True}{\lnot !B} \in \Gamma$. La branche prolongée contient
  alors les !!{signed formula}s de
  $\Gamma \cup \{\sFmla{\False}{!B}\}$. Supposons
  $\iftag{FOL}{\Sat{M}}{\pSat{v}}{\Gamma}$. En particulier,
  $\iftag{FOL}{\Sat{M}}{\pSat{v}}{\lnot !B}$. Par conséquent,
  $\iftag{FOL}{\Sat/{M}}{\pSat/{v}}{!B}$, c'est-à-dire que
  $\iftag{FOL}{\Struct{M}}{\pAssign{v}}$ satisfait
  $\sFmla{\False}{!B}$.
\item La branche est prolongée par l'application de
  $\TRule{\False}{\lnot}$ à
  $\sFmla{\False}{\lnot !B} \in \Gamma$ : exercice.
\item La branche est prolongée par l'application de
  $\TRule{\True}{\land}$ à
  $\sFmla{\True}{!B \land !C} \in \Gamma$. Deux nouvelles
  !!{signed formula}s apparaissent sur la branche :
  $\sFmla{\True}{!B}$ et $\sFmla{\True}{!C}$. Supposons
  $\iftag{FOL}{\Sat{M}}{\pSat{v}}{\Gamma}$, et en particulier
  $\iftag{FOL}{\Sat{M}}{\pSat{v}}{!B \land !C}$. Alors
  $\iftag{FOL}{\Sat{M}}{\pSat{v}}{!B}$ et
  $\iftag{FOL}{\Sat{M}}{\pSat{v}}{!C}$. Ainsi,
  $\iftag{FOL}{\Struct{M}}{\pAssign{v}}$ satisfait à la fois
  $\sFmla{\True}{!B}$ et $\sFmla{\True}{!C}$.
\item La branche est prolongée par l'application de
  $\TRule{\False}{\lor}$ à
  $\sFmla{\False}{!B \lor !C} \in \Gamma$ : exercice.
\item La branche est prolongée par l'application de
  $\TRule{\False}{\lif}$ à
  $\sFmla{\False}{!B \lif !C} \in \Gamma$. Deux nouvelles
  !!{signed formula}s apparaissent sur la branche :
  $\sFmla{\True}{!B}$ et $\sFmla{\False}{!C}$. Supposons
  $\iftag{FOL}{\Sat{M}}{\pSat{v}}{\Gamma}$, et en particulier
  $\iftag{FOL}{\Sat/{M}}{\pSat/{v}}{!B \lif !C}$. Alors
  $\iftag{FOL}{\Sat{M}}{\pSat{v}}{!B}$ et
  $\iftag{FOL}{\Sat/{M}}{\pSat/{v}}{!C}$. Ainsi,
  $\iftag{FOL}{\Struct{M}}{\pAssign{v}}$ satisfait à la fois
  $\sFmla{\True}{!B}$ et $\sFmla{\False}{!C}$.

\iftag{FOL}{%
\item La branche est prolongée par l'application de
  $\TRule{\True}{\lforall}$ à
  $\sFmla{\True}{\lforall[x][!A(x)]} \in \Gamma$. La nouvelle
  !!{signed formula}~$\sFmla{\True}{!A(t)}$ est ajoutée à la branche.
  Supposons $\Sat{M}{\Gamma}$, et en particulier
  $\Sat{M}{\lforall[x][!A(x)]}$. D'après
  \olref[syn][sem]{prop:quant-terms}, $\Sat{M}{!A(t)}$. Par conséquent,
  $\Struct{M}$ satisfait $\sFmla{\True}{!A(t)}$.

\item La branche est prolongée par l'application de
  $\TRule{\False}{\lforall}$ à
  $\sFmla{\False}{\lforall[x][!A(x)]} \in \Gamma$. On ajoute la
  !!{signed formula}~$\sFmla{\False}{!A(a)}$, où $a$ est une
  !!{constant} qui ne figure pas dans~$\Gamma$. Comme $\Gamma$ est
  satisfaisable, il existe une $\Struct{M}$ telle que
  $\Sat{M}{\Gamma}$, et en particulier
  $\Sat/{M}{\lforall[x][!A(x)]}$. Il faut montrer que
  $\Gamma \cup \{\sFmla{\False}{!A(a)}\}$ est satisfaisable. Pour ce
  faire, définissons une structure appropriée~$\Struct{M'}$.

  D'après \olref[syn][ass]{prop:sat-quant},
  $\Sat/{M}{\lforall[x][!A(x)]}$ si et seulement s'il existe une
  assignation~$s$ telle que $\Sat/{M}{!A(x)}[s]$. Soit maintenant
  $\Struct{M'}$ la structure identique à $\Struct{M}$, sauf que
  $\Assign{a}{M'} = s(x)$. D'après
  \olref[syn][ext]{cor:extensionality-sent}, pour toute
  $\sFmla{\True}{!C} \in \Gamma$, on a $\Sat{M'}{!C}$, et pour toute
  $\sFmla{\False}{!C} \in \Gamma$, on a $\Sat/{M'}{!C}$, puisque $a$
  ne figure pas dans~$\Gamma$.

  D'après \olref[syn][ext]{prop:extensionality},
  $\Sat/{M'}{!A(x)}[s]$. Puis
  \olref[syn][ext]{prop:ext-formulas} donne
  $\Sat/{M'}{!A(a)}[s]$. Comme $!A(a)$ est !!a{sentence},
  \olref[syn][ass]{prop:sentence-sat-true} entraîne
  $\Sat/{M'}{!A(a)}$. Autrement dit, $\Struct{M'}$ satisfait
  $\sFmla{\False}{!A(a)}$.
\item La branche est prolongée par l'application de
  $\TRule{\True}{\lexists}$ à
  $\sFmla{\True}{\lexists[x][!B(x)]} \in \Gamma$ : exercice.
\item La branche est prolongée par l'application de
  $\TRule{\False}{\lexists}$ à
  $\sFmla{\False}{\lexists[x][!B(x)]} \in \Gamma$ : exercice.}{}
\end{enumerate}
Considérons maintenant les inférences qui provoquent une ramification.
\begin{enumerate}
\item La branche est prolongée par l'application de
  $\TRule{\False}{\land}$ à
  $\sFmla{\False}{!B \land !C} \in \Gamma$. Elle se ramifie en une
  branche gauche qui se poursuit par $\sFmla{\False}{!B}$ et une
  branche droite qui se poursuit par $\sFmla{\False}{!C}$. Supposons
  $\iftag{FOL}{\Sat{M}}{\pSat{v}}{\Gamma}$, et en particulier
  $\iftag{FOL}{\Sat/{M}}{\pSat/{v}}{!B \land !C}$. Alors
  $\iftag{FOL}{\Sat/{M}}{\pSat/{v}}{!B}$ ou
  $\iftag{FOL}{\Sat/{M}}{\pSat/{v}}{!C}$. Dans le premier cas,
  $\iftag{FOL}{\Struct{M}}{\pAssign{v}}$ satisfait
  $\sFmla{\False}{!B}$, donc toutes les formules de la branche gauche.
  Dans le second,
  $\iftag{FOL}{\Struct{M}}{\pAssign{v}}$ satisfait
  $\sFmla{\False}{!C}$, donc toutes les formules de la branche droite.
\item La branche est prolongée par l'application de
  $\TRule{\True}{\lor}$ à
  $\sFmla{\True}{!B \lor !C} \in \Gamma$ : exercice.
\item La branche est prolongée par l'application de
  $\TRule{\True}{\lif}$ à
  $\sFmla{\True}{!B \lif !C} \in \Gamma$ : exercice.
\item La branche est prolongée au moyen de \Cut{}. On obtient deux
  branches, dont l'une contient $\sFmla{\True}{!B}$ et l'autre
  $\sFmla{\False}{!B}$. Comme
  $\iftag{FOL}{\Sat{M}}{\pSat{v}}{\Gamma}$ et que l'on a soit
  $\iftag{FOL}{\Sat{M}}{\pSat{v}}{!B}$, soit
  $\iftag{FOL}{\Sat/{M}}{\pSat/{v}}{!B}$,
  $\iftag{FOL}{\Struct{M}}{\pAssign{v}}$ satisfait la branche gauche
  ou la branche droite.
\end{enumerate}
\end{proof}

\begin{editorial}
Dans les deux cas consacrés à $\lforall$, la source commence avec
$!B(x)$, puis poursuit le même argument avec $!A(x)$ et ses instances.
La lettre~$A$ est rétablie dans les prémisses afin que chaque règle et
chaque calcul sémantique portent sur une seule formule. Dans le passage
commun aux présentations propositionnelle et du premier ordre,
« structure » est adapté en « valuation » pour la première.
\end{editorial}

\tagprob{FOL}
\begin{prob}
Achever la démonstration de
\olref[fol][tab][sou]{thm:tableau-soundness}.
\end{prob}
\tagendprob

\tagprob{notFOL}
\begin{prob}
Achever la démonstration de
\olref[pl][tab][sou]{thm:tableau-soundness}.
\end{prob}
\tagendprob

\begin{cor}
\ollabel{cor:weak-soundness}
Si $\Proves !A$, alors $!A$ est \iftag{FOL}{valide}{une tautologie}.
\end{cor}

\begin{cor}
\ollabel{cor:entailment-soundness}
Si $\Gamma \Proves !A$, alors $\Gamma \Entails !A$.
\end{cor}

\begin{proof}
  Si $\Gamma \Proves !A$, il existe $!B_1$, \dots, $!B_n \in
  \Gamma$ tels que $\{\sFmla{\False}{!A},
  \sFmla{\True}{!B_1}, \dots, \sFmla{\True}{!B_n}\}$ possède un
  !!{tableau} fermé. D'après
  \olref{thm:tableau-soundness}, toute
  \iftag{FOL}{!!{structure}}{!!{valuation}}~$\iftag{FOL}{\Struct{M}}{\pAssign{v}}$
  rend faux au moins un $!B_i$, ou rend vrai $!A$. Par conséquent, si
  $\iftag{FOL}{\Sat{M}}{\pSat{v}}{\Gamma}$, alors
  $\iftag{FOL}{\Sat{M}}{\pSat{v}}{!A}$.
\end{proof}

\begin{cor}
\ollabel{cor:consistency-soundness}
Si $\Gamma$ est satisfaisable, alors il est cohérent.
\end{cor}

\begin{proof}
Nous démontrons la contraposée. Supposons que $\Gamma$ soit
incohérent. Il existe alors $!B_1$, \dots, $!B_n \in \Gamma$ et un
!!{tableau} fermé pour
$\{\sFmla{\True}{!B_1}, \dots, \sFmla{\True}{!B_n}\}$. D'après
\olref{thm:tableau-soundness}, il n'existe aucune
$\iftag{FOL}{\Struct{M}}{\pAssign{v}}$ telle que
$\iftag{FOL}{\Sat{M}}{\pSat{v}}{!B_i}$ pour tout
$i=1$, \dots,~$n$. L'ensemble $\Gamma$ n'est donc pas satisfaisable.
\end{proof}

\end{document}
